\section{Branch-and-Benders and graph-based inequalities}
\label{sec:methods}

The extensive model creates burn variables for every sampled cell and
scenario. We instead retain the common placement in a master problem and
estimate each scenario's recourse loss. Structural graph inequalities then
raise those estimates before the full scenario LP has been processed.
The graph results in this section apply to a finite directed graph unless
the statement explicitly requires a rooted tree.

\subsection{Exact decomposition}\label{sec:benders}

Introduce a nonnegative master variable $\eta_s$ for the loss of scenario
$s$. The remaining variables $t,z_s$ aggregate risk exactly as in
\eqref{eq:extensive}. In particular, the master objective and risk
constraints are
\begin{equation}\label{eq:master}
 \begin{aligned}
 \min_{y\in\mathcal Y,\eta,t,z}\quad
 &(1-\lambda)\sum_s\rho_s\eta_s
   +\lambda\left(t+\frac{1}{1-\beta}\sum_s\rho_sz_s\right)\\
 \text{s.t.}\quad
 &z_s\ge\eta_s-t,\quad z_s\ge0,\quad\eta_s\ge0
       &&(s\in\mathcal S),\qquad t\in\mathbb R.
 \end{aligned}
\end{equation}
The missing constraints $\eta_s\ge Q_s(y)$ are supplied by Benders
optimality cuts. If $\bar y\in[0,1]^C$, denote by $\overline Q_s(\bar y)$
the value of the propagation LP \eqref{eq:propagation} with $y=\bar y$.
The LP dual generates affine global underestimators
\begin{equation}\label{eq:benders-cut}
 \eta_s\ge \overline Q_s(\bar y)
         +\sum_{i\in C}\pi_{si}(y_i-\bar y_i),
\end{equation}
where $\pi_s$ is the vector of multipliers of the rows fixing the
subproblem's copies of $y$ to $\bar y$. The cut is tight at $\bar y$ and
valid for all feasible placements. The implementation solves scenario LPs
and adds violated cuts at integer candidate solutions through a callback;
configured variants also add LP-dual user cuts at the root relaxation.
Proposition~\ref{prop:recourse} makes the integer-candidate checks equal
to the actual scenario losses. Thus, with exact subproblem solves and the
stated solver tolerances, the candidate cuts support an exact
branch-and-cut solution of the finite-scenario problem.

The inequalities below concern individual scenario estimates $\eta_s$.
They remain valid under expected loss, CVaR and mean--CVaR because those
models change only the master risk aggregation. They strengthen the
implemented search but do not replace its integer-candidate Benders
checks. The theoretical equality of a \emph{complete} inequality family
with recourse must not be confused with adding a limited number of root
cuts \citep{rahmaniani2017benders}.

\subsection{Downstream and capped-coverage bounds}
\label{sec:coverage}

Let $K_s=R_s(0)$ be the untreated reachable cells and
$q_s^0=\sum_{k\in K_s}w_k$. For eligible $i\in C\cap
(K_s\setminus\{r_s\})$, define its \emph{closed downstream set}
$D_s(i)=\{k\in K_s:\text{there is a directed path from $i$ to $k$}
\text{ in }G_s\}$. In particular $i\in D_s(i)$. For other eligible cells,
put $D_s(i)=\varnothing$, and set
$d_{si}=\sum_{k\in D_s(i)}w_k$. These are structural potential
protection masses, not necessarily the actual losses avoided by one break
in a graph with several routes.

The \emph{standard lifted lower-bound inequality} (LLBI) is
\begin{equation}\label{eq:standard}
 \eta_s\ge L_s^{\rm std}(y):=
           q_s^0-\sum_{i\in C}d_{si}y_i.
\end{equation}
Its construction is cheap once downstream sets have been computed, but
several coefficients may count the same potentially protected cell.

\begin{proposition}[Validity of the standard LLBI]
\label{prop:standard}
For every binary placement $y$, inequality \eqref{eq:standard} is valid.
It is also valid below the fractional recourse function
$\overline Q_s(y)$ for $y\in[0,1]^C$.
\end{proposition}
\begin{proof}
If a cell $k\in K_s$ is no longer reachable under a binary placement,
each original ignition-to-$k$ path meets a selected non-ignition cell.
At least one such cell $i$ has $k\in D_s(i)$. Hence the protected cells
lie in $\bigcup_{i:y_i=1}D_s(i)$, whose weight is at most
$\sum_i d_{si}y_i$. This proves the first statement.
For fractional $y$, sum the propagation constraints along any original
ignition-to-$k$ path $p$ to obtain
$x_{sk}\ge 1-\sum_{i\in C(p)}y_i$, where $C(p)$ is the eligible
non-ignition part of the path. Every $i\in C(p)$ belongs to the set of
candidates downstream of which $k$ lies. Therefore
$x_{sk}\ge 1-\sum_{i\in A_{sk}}y_i$, where $A_{sk}$ is defined below.
Summing the nonnegative-weight inequalities over $K_s$ yields
\eqref{eq:standard} after interchanging the two sums.
\end{proof}

To avoid repeatedly counting potential protection of a cell, define
\begin{equation}\label{eq:coverage-sets}
 A_{sk}=\{i\in C\cap(K_s\setminus\{r_s\}): k\in D_s(i)\},
 \qquad K_s^+=\{k\in K_s:A_{sk}\ne\varnothing\}.
\end{equation}
On a graph with alternative paths, $i\in A_{sk}$ means that $i$ \emph{can}
lie upstream of $k$; it does not mean that treating $i$ alone prevents
every route to $k$. For each $k\in K_s^+$, introduce
$0\le\zeta_{sk}\le1$ and impose
\begin{equation}\label{eq:coverage-ext}
 \eta_s\ge q_s^0-\sum_{k\in K_s^+}w_k\zeta_{sk},
 \qquad
 \zeta_{sk}\le\sum_{i\in A_{sk}}y_i\quad(k\in K_s^+).
\end{equation}
The cap at one prevents the same cell from receiving more than its own
weight in the bound.

\begin{proposition}[Coverage projection and separation]
\label{prop:coverage}
For every $y\in[0,1]^C$, the strongest bound implied by
\eqref{eq:coverage-ext} is
\begin{equation}\label{eq:coverage-envelope}
 L_s^{\rm cov}(y)=q_s^0-
  \sum_{k\in K_s^+}w_k
       \min\left\{1,\sum_{i\in A_{sk}}y_i\right\}
 \le\overline Q_s(y).
\end{equation}
Its exact projection is the linear family
\begin{equation}\label{eq:coverage-proj}
 \eta_s\ge q_s^0-
   \sum_{k\in K_s^+\setminus U}w_k
   -\sum_{k\in U}w_k\sum_{i\in A_{sk}}y_i,
 \quad U\subseteq K_s^+.
\end{equation}
At $\bar y$, a strongest member takes
$U=\{k:\sum_{i\in A_{sk}}\bar y_i<1\}$, with either choice allowed at
equality. Taking $U=K_s^+$ gives exactly \eqref{eq:standard}.
\end{proposition}
\begin{proof}
The path argument in Proposition~\ref{prop:standard}, together with
$x_{sk}\ge0$, gives
$x_{sk}\ge[1-\sum_{i\in A_{sk}}y_i]^+$ for $k\in K_s^+$;
for $k\notin K_s^+$ the same path argument gives $x_{sk}=1$.
This proves the inequality in \eqref{eq:coverage-envelope}.
For fixed $y$, maximizing each $\zeta_{sk}$ subject to
\eqref{eq:coverage-ext} permits
$\zeta_{sk}=\min\{1,\sum_{i\in A_{sk}}y_i\}$.
Since $q_s^0-w_k\min\{1,a_k\}=
\max\{q_s^0-w_k,q_s^0-w_ka_k\}$ cellwise for $w_k\ge0$,
independent selections of the active term give the maximum of the
right-hand sides of \eqref{eq:coverage-proj}. Finally,
$\sum_{k\in K_s^+}w_k\sum_{i\in A_{sk}}y_i
=\sum_i\sum_{k\in D_s(i)}w_ky_i=\sum_i d_{si}y_i$.
\end{proof}

We call the single static member $U=K_s^+$ CoverageLLBI-poly, and
separation of the full family at root LP points CoverageLLBI-exp.
The static projected member and the standard LLBI are identical when
defined on the same scenario and candidate set. Their complete projected
family can nevertheless improve the master bound through cellwise caps.

\subsection{Propagation-path bounds}\label{sec:path}

Coverage uses all candidates that might reach a cell. A complementary
construction uses the actual routes along which a fire \emph{can} reach
it. Let $\mathcal P_{sk}$ be all directed simple paths from $r_s$ to
$k\in K_s$ and $\widehat{\mathcal P}_{sk}\subseteq\mathcal P_{sk}$
the paths retained by an implementation. Put
$C(p)=C\cap(V(p)\setminus\{r_s\})$; thus an eligible terminal cell
belongs to $C(p)$, while the ignition never does. The trivial path to
$r_s$ has $C(p)=\varnothing$. A path inequality
\begin{equation}\label{eq:path-atom}
 b_{sk}+\sum_{i\in C(p)}y_i\ge1,
 \qquad p\in\widehat{\mathcal P}_{sk},\quad 0\le b_{sk}\le1,
\end{equation}
forces a lower estimate $b_{sk}$ of burning whenever that path is
unblocked. Coupling the estimates by
$\eta_s\ge\sum_{k\in K_s}w_kb_{sk}$ gives the extended PathLLBI.

To describe the projection, allow the symbol $\bot$ to omit a cell from a
cut, and let $\pi_k\in\widehat{\mathcal P}_{sk}\cup\{\bot\}$ be an
independent choice. If $K(\pi)=\{k:\pi_k\ne\bot\}$, the resulting
projected family is
\begin{equation}\label{eq:path-proj}
 \eta_s\ge\sum_{k\in K(\pi)}w_k
       \left(1-\sum_{i\in C(\pi_k)}y_i\right)
 \quad\left(\pi\in\prod_{k\in K_s}
        (\widehat{\mathcal P}_{sk}\cup\{\bot\})\right).
\end{equation}

\begin{proposition}[Path projection and complete-path recourse]
\label{prop:path}
The extended PathLLBI and \eqref{eq:path-proj} define the same
$(y,\eta_s)$ projection. Every member is valid below
$\overline Q_s(y)$ for $y\in[0,1]^C$. If every simple path is retained,
then this projection is exactly the epigraph of
$\overline Q_s(y)$.
\end{proposition}
\begin{proof}
Summing \eqref{eq:arcs} along a path proves
$x_{sk}\ge1-\sum_{i\in C(p)}y_i$; $x_{sk}\ge0$ supplies the
$\bot$ term. With nonnegative weights, the minimum feasible $b_{sk}$ is
\[
 \max\left\{0,\max_{p\in\widehat{\mathcal P}_{sk}}
       \left(1-\sum_{i\in C(p)}y_i\right)\right\}.
\]
Maximizing independently over these terms for each $k$ yields
\eqref{eq:path-proj}, and also proves validity. For completeness, with
all paths retained define $d_k(y)=\min_{p\in\mathcal P_{sk}}
\sum_{i\in C(p)}y_i$ and $x_{sk}^*=[1-d_k(y)]^+$.
The trivial path gives $x_{sr_s}^*=1$. For each arc $(i,j)$,
$d_j(y)\le d_i(y)+\bar y_{sj}$: appending the arc to a minimum-cost
ignition-to-$i$ route gives a walk, and removing any nonnegative-cost
cycles cannot increase its cost. Hence $x_{si}^*\le
x_{sj}^*+\bar y_{sj}$; the bound also holds for arcs into the fixed
ignition and arcs whose tail is unreachable. Thus $x^*$ satisfies the
recourse LP and attains every cell's individual path lower bound.
Summing with nonnegative weights proves equality with its value.
\end{proof}

Complete PathLLBI-exp would choose a minimum-cost path for each cell with
$d_k(\bar y)<1$ when separating at a fractional point $\bar y$.
The implemented projected separator instead selects a least-cost path
among the paths it has \emph{stored} for each destination. It first
enumerates up to a configured number of simple paths per cell, so its
separation need not attain the complete-path envelope on a general
graph. PathLLBI-poly keeps one deterministic first path per destination
and one static projected cut. Both versions are valid because omitted
paths weaken the bound; they do not create an invalid cut.

\subsection{Bound relationships and the tree case}\label{sec:relationships}

The relations explain what strengthening is mathematically possible and
what a computation can instead reveal about implementation overhead.
For $y\in[0,1]^C$, let $L_s^{\rm path}(y)$ denote the complete-path
envelope. Then
\begin{equation}\label{eq:hierarchy}
 L_s^{\rm std}(y)\le L_s^{\rm cov}(y)
          \le L_s^{\rm path}(y)=\overline Q_s(y).
\end{equation}
The first inequality follows since
$\min\{1,a\}\le a$. For each path $p$ to $k$, its candidate set
$C(p)$ is contained in $A_{sk}$; therefore
$[1-\sum_{i\in A_{sk}}y_i]^+
\le[1-\sum_{i\in C(p)}y_i]^+$. Summing cellwise with nonnegative
weights proves the second inequality. Proposition~\ref{prop:path}
gives equality at the right. These comparisons concern envelopes;
a particular finite collection of cuts can produce weaker master bounds.

On a rooted propagation tree each cell has one ignition path $p_{sk}$,
and precisely its eligible non-ignition ancestors reach $k$. Thus
$A_{sk}=C(p_{sk})$ and
\begin{equation}\label{eq:tree-equality}
 L_s^{\rm cov}(y)=L_s^{\rm path}(y)=\overline Q_s(y)
 \qquad(y\in[0,1]^C).
\end{equation}
Moreover, the all-unsaturated coverage cut and the one-path-per-cell
static path cut both equal the standard LLBI. Consequently, a runtime
difference between separated CoverageLLBI and PathLLBI on a \emph{verified}
tree cannot by itself demonstrate stronger \emph{complete} bounds. Finite
separation, cut selection, construction time and solver interaction are
distinct sources of computational differences. The actual reduced graphs
are predominantly DAGs, on which complete PathLLBI can be strictly
stronger than complete CoverageLLBI. The implemented path family stores
a limited number of paths per cell; its separation need not attain the
complete-path envelope.
